% Transcription — fonds Alexandre Grothendieck, shelfmark 156-5, batch 2 (pages 21-40).
% Established from the facsimile archives/batches/156-5/batch-02.pdf (a mirror of
% https://grothendieck.umontpellier.fr/156-5.pdf; no cover sheet in the batch, so
% PDF page k = folder page 20+k), per the skill transcribe-grothendieck.
%
% All twenty leaves are mathematics and all are transcribed; there is no cover
% and no unrelated verso. His own pagination 21-40 runs across the top of the
% leaves and coincides with the archivists'. One dated heading: « 18 juin »
% in the margin of page 36. Page 21 picks up mid-proof where batch 1 broke off
% (b) ⇒ a) of the proposition on raffinement, page 20).
%
% Content. Pages 21-22: end of the proof that the two forms a), b) of
% « F raffine G » agree, a corollary b') (∀A ∃B, A° ⊂ B°, and A° ⊂ B ⇒ A ⊂ B),
% the definition of raffinement F ≪ G (an order on figures) and a first
% proposition on maximal raffinements. Pages 23-26: subdivision F ≼ G
% (F ≪ G and |F| = |G|, equivalently maximal among raffinements), the
% dictionary with a diagram S ⊂ T ⊂ L over an increasing φ : I → J; every
% raffinement is a sub-figure of a subdivision (F ≤ G' ≼ G, G' obtained by
% adding the singletons of |G| ∖ |F|); ≪ and ≼ determine one another through ≤,
% and F ≤ G, F ≼ G ⇒ F = G; the axioms C5-C8 checked. Pages 27-29: an attempted
% description of |F| ∩ |G| (struck), then the « figure intersection » F.G, an
% Inf for ≪ (not in general F ∩ G), stable under ≪, ≼, ≤ in one variable;
% corollary: axiom C6 holds; statement of C7 and of an auxiliary C7₀.
% Pages 29-33: for G ≤ F and a subdivision G' ≼ G, the subdivision
% F' = G' ∪ (F ∖ G) of F inducing G' (proved to be a figure and a
% subdivision), then a corollary characterising it by four equivalent
% conditions (a maximality, b bijection F' ∖ G' ≅ F ∖ G, c strata, d F' ⊃ F ∖ G),
% d ⇒ b and b ⇔ c ⇔ d proved; a ⇔ the rest begun and struck (p. 33). Pages
% 34-36: the same in terms of factorisations S → I' → I and T → J' → J, the
% ordered set K = J' ⊔ (I ∖ J) as final object. Pages 36-40 (18 June): figures
% recast as ultra-topologies (Alexandrov topologies) (S, T) on parts S of L;
% sub-figures = closed parts, ≪ = continuous inclusion, ≼ = finer topology
% on the same set, Inf = initial topology; for T ⊂ S closed and a finer
% topology T'₀ on T, the coarsest topology on S finer than the given one and
% inducing T'₀, with its closed sets described by three conditions (p. 40,
% where the batch stops mid-statement at c)).
%
% Reading hazards fixed once here. F, G, H, F', G' are figures; A° the stratum
% of A; |F| the support; R_F the equivalence relation of strata. ≪ is
% raffinement, ≼ (his ≤ with a curl, also drawn ≺ over =) is subdivision, ≤ is
% sub-figure. From page 38 on he writes 𝒯 for a topology while T is the set
% |G|; on pages 36-37 the topology of a figure (S, T) is still written T.
% 𝔗 on page 40 renders a looped, underlined 𝒯 for the set of topologies. The
% asterisk-like mark on A (pp. 21, 22, 30) is kept as A^*; it may be a
% heavily inked °. His underlining is set in emphasis. The hand is fast
% drafting throughout; slanted left-margin notes on pages 21-23, 25, 29,
% 31, 33, 35, 38 (a vertical note over the full height) and 39 were turned
% and read in part.
%
% Readings flagged and not settled: most of the page-22 and page-23 margins;
% the second sentence of the corollary on page 23 (« compatibles », « sup. »);
% the letter 𝒥 for the set of figures (p. 24); the struck end of page 35;
% the order of the interlinear additions at the foot of page 38; the reading
% « demi-treillis » (p. 28). Variances on the page and not repaired: I × J
% where I × I' is meant (p. 27); in the corollary of page 31, a) is crossed
% through but still used as a) in the proof plan.
%
% Pass: Opus 5.5 (claude-opus-5-5), 2026-09-24 — first pass, unchecked against the pages by a human.
\documentclass[11pt,a4paper]{article}
% Le préambule commun à toutes les transcriptions du fonds Grothendieck.
%
% Il tient en une idée : **l'incertitude se compose, elle ne se résout pas.**
% Une transcription qui lisse un mot douteux a détruit la seule chose pour
% laquelle on la faisait — un lecteur doit pouvoir savoir, à chaque endroit,
% ce qui était sur le feuillet et ce qui a été ajouté. Les six macros qui
% suivent sont donc l'appareil critique, et non de la décoration.
%
% Les mêmes commandes sont reconnues par `scripts/render.mjs`, qui produit la
% vue de lecture du site. Une macro ajoutée ici doit l'être aussi là-bas,
% faute de quoi le rendu échouera bruyamment — ce qui est voulu : un rendu
% silencieusement faux est pire qu'un rendu absent.

\ProvidesPackage{grothendieck}[2026/08/08 Transcriptions du fonds Grothendieck]

\usepackage{amsmath,amssymb,amsthm}
\usepackage{mathrsfs}
\usepackage{stmaryrd}
% Les diagrammes commutatifs sont partout dans ce fonds : c'est la langue
% même de la géométrie algébrique de Grothendieck, et une transcription qui
% les rendrait en prose ne transcrirait rien.
\usepackage{tikz-cd}
% Français par défaut — c'est la langue des feuillets — mais l'anglais est
% chargé aussi : la traduction et le résumé sont des documents anglais, et la
% typographie française (espaces avant les deux-points) y serait une faute.
% Ces documents basculent par \selectlanguage{english} en tête de corps.
\usepackage[english,french]{babel}
\usepackage{fontspec}
\usepackage{geometry}
\geometry{a4paper,margin=25mm,marginparwidth=28mm}
\usepackage{marginnote}
\usepackage{xcolor}
% ulem, not soul. `soul`'s \st and \ul are built on a character-by-character
% scan that XeLaTeX's font machinery breaks: the document fails with an
% undefined control sequence rather than degrading. ulem does the same two jobs
% and is Unicode-engine safe. `normalem` stops it hijacking \emph, which the
% transcriptions use for Grothendieck's own emphasis.
\usepackage[normalem]{ulem}
\usepackage{hyperref}
\hypersetup{colorlinks=true,linkcolor=black,urlcolor=black,pdfborder={0 0 0}}

% --- Métadonnées du lot -----------------------------------------------------
% Reprises telles quelles par le script de rendu, qui les affiche en tête de
% la vue de lecture. Elles fixent ce que le fichier prétend couvrir : sans
% elles, une transcription flotte sans point d'attache dans un fonds de seize
% mille feuillets.

\newcommand{\folder}[1]{\gdef\grothFolder{#1}}
\newcommand{\batch}[1]{\gdef\grothBatch{#1}}
\newcommand{\leaves}[2]{\gdef\grothFirst{#1}\gdef\grothLast{#2}}
\newcommand{\foldertitle}[1]{\gdef\grothFoldertitle{#1}}
\gdef\grothFolder{?}\gdef\grothBatch{?}\gdef\grothFirst{?}\gdef\grothLast{?}\gdef\grothFoldertitle{}

% --- Le résumé --------------------------------------------------------------
%
% Il ouvre la lecture modernisée, sous le titre « Résumé », et il est composé
% en petit. Ce n'est pas un ornement : c'est le seul passage du document qui
% ne soit pas des mathématiques, et un lecteur doit pouvoir le reconnaître
% avant de l'avoir lu — puis le sauter s'il connaît déjà le sujet, ou s'y
% arrêter s'il ne le connaît pas. Le filet qui le suit marque la reprise du
% corps.

\newenvironment{resume}
  {\small\setlength{\parskip}{0.45em}}
  {\par\medskip\hrule\medskip}

% --- Mots-clés ---------------------------------------------------------------
%
% En anglais, en clôture du résumé de la lecture modernisée : le vocabulaire
% moderne sous lequel chercher ce que les feuillets construisent. Le manifeste
% du site (`npm run manifest`) les extrait pour étiqueter la cote entière —
% cette ligne est la source unique des tags, il n'y a pas de fichier à côté.
\newcommand{\keywords}[1]{\par\smallskip\noindent
  {\footnotesize\textsc{Keywords} --- \textit{#1}}\par}

% --- L'appareil critique ----------------------------------------------------

% Le numéro de feuillet, en marge. C'est l'ancre qui permet de régler un
% désaccord sur une formule en regardant *un* feuillet plutôt qu'une chemise
% de six cents. Le site s'en sert aussi pour tourner les pages du fac-similé
% quand on fait défiler la transcription.
\newcommand{\leaf}[1]{%
  \marginnote{\footnotesize\color{gray!70}\textsf{#1}}%
  \label{leaf:#1}\ignorespaces}

% Illisible : on ne devine pas. Un mot inventé qui se lit comme les autres est
% le pire résultat possible.
\newcommand{\ill}{\textcolor{red!60!black}{[\,\ldots\,]}}

% Lecture douteuse : le mot est donné, et signalé comme incertain.
\newcommand{\uncertain}[1]{\uline{#1}}

% Ajout de l'éditeur — une lettre manquante, un mot sous-entendu.
\newcommand{\add}[1]{\textcolor{blue!50!black}{[#1]}}

% Biffé par Grothendieck lui-même. Ce qu'il a rayé fait partie du document :
% une rature dit souvent où la pensée a changé de direction.
\newcommand{\struck}[1]{\sout{#1}}

% Note du transcripteur, sur l'état matériel du feuillet.
\newcommand{\note}[1]{\par\smallskip{\small\color{gray!60!black}\textit{[#1]}}\par\smallskip}

% Note marginale de Grothendieck, distincte du corps du texte.
\newcommand{\marginal}[1]{\par\smallskip\noindent\hspace{1em}%
  {\small\color{gray!60!black}#1}\par\smallskip}

% --- Titre ------------------------------------------------------------------

\AtBeginDocument{%
  \begin{center}
    {\large\bfseries Fonds Alexandre Grothendieck --- cote n\textsuperscript{o}~\grothFolder}\\[2pt]
    {\itshape\grothFoldertitle}\\[4pt]
    {\small feuillets \grothFirst--\grothLast\ (lot \grothBatch)}\\[2pt]
    {\footnotesize\color{gray!60!black}Transcription établie d'après le fac-similé numérisé
     par l'Université de Montpellier.\\
     Les lectures incertaines sont soulignées, les illisibles notées [\,\ldots\,],
     les ajouts entre crochets.}
  \end{center}
  \vspace{1em}}

\endinput


\folder{156-5}
\batch{2}
\pages{21}{40}
\dating{1986}
\watermark{Édition de démonstration}
\foldertitle{[Chapitre] V. Algèbre des figures : notes manuscrites (14/06/1986).}

\begin{document}
\page{21}
On \uncertain{sait} $B' \in G$ est tel que $A'^{\circ} \subset B'^{\circ}$, donc \struck{$A \subset B'$} (hyp. dans b)) $B' \subset B$, donc $A'^{\circ} \subset B$, OK.

Prouvons que $B$ est le plus petit él. de $G^{A}$. Soit $B' \in G^{A}$ i.e. $B' \in G$, $A \subset B'$, prouvons $B \subset B'$ \uncertain{ou encore} $B^{\circ} \cap B' \neq \emptyset$ \struck{Comme}. Comme $A^{\circ} \subset B^{\circ}$ et $A^{\circ} \subset A \subset B'$, on a $\underset{\neq \emptyset}{A^{\circ}} \subset B^{\circ} \cap B'$, OK.

Prouvons ensuite que $A \in F$, $B \in G$ implique que $A \cap B$ est $F$-fermé dans $|F|$. Par b) on voit que $\underset{\displaystyle \bigcup_{A \in F} A^{\circ}}{\overset{\|}{|F|}} \subset |G|$, et que la relation d'équiv. $R_{F}$ sur $|F|$ est plus fine que celle induite par $R_{G}$ sur $|G|$. Cela implique que $A \cap B$ est saturé pour la relation d'équiv. $R_{F}$ (puisque $B$ l'est pour $R_{G}$). Donc
\[
  A \cap B = \bigcup_{\substack{C \in F \\ \text{t.q. } C^{\circ} \subset A \cap B}} C^{\circ} .
\]
Je dis que \struck{\ill{}} dans cette réunion, \struck{\ill{}} on peut remplacer $C^{\circ}$ par $C$ i.e.
\[
  C \in F, \quad C^{\circ} \subset A \cap B \Longrightarrow C \subset A \cap B .
\]
$C \subset A$\note{l'exposant de ce $C$ paraît biffé}. Reste : prouver \struck{\ill{}} \ill{} $C^{\circ} \subset A$, \ill{} $C^{\circ} \subset B$. On en déduit un complément à l'énoncé de la prop. :

\textbf{Corollaire} La condition b) équivaut à
\begin{itemize}
  \item[b')] $\forall A \in F$ $\exists B \in G$ tel que $A^{\circ} \subset B^{\circ}$, et d'autre part $\forall A \in F$, $B \in G$, $A^{\circ} \subset B$ implique $A \subset B$.
\end{itemize}
\marginal{On a donc : (b'') $\Longleftrightarrow$ : $|F| \subset |G|$, la relation de \uncertain{préordre} sur $|F|$ définie par $F$ est plus fine que celle induite \ill{} relation \ill{} $|G|$}
\note{note oblique dans le coin inférieur gauche, qui chevauche la fin des lignes du texte}

\page{22}
Il nous fallait prouver b $\Rightarrow$ b' pour achever la dém. de la prop. \struck{Voir \ill{} \ill{} \ill{}} \struck{\ill{} voir \ill{} (prouvons b $\Rightarrow$ b') que, supposant b), $A \in F$, $B \in G$, $A^{\circ} \subset B \Longrightarrow A \subset B$ (cf} On l'a prouvé \struck{\ill{}} dans la dém. de b) $\Rightarrow$ a) au cours (p. 20)\note{un trait encadre « supposant b) … $A \subset B$ » et le renvoie après « On l'a prouvé »}. Reste : prouver b' $\Rightarrow$ b donc, supposant b', le fait que $\varphi$ est croissant :
\[
  A, A' \in F, \quad B, B' \in G, \quad A^{\circ} \subset B^{\circ}, \quad A'^{\circ} \subset B'^{\circ}, \quad A' \subset A^{*} \Longrightarrow B' \subset B^{*}
\]
\note{les astérisques sont de sa main ; une seconde ligne, en partie biffée, finit par « $A^{*} \subset B^{*}$ »}
\struck{\ill{} b') \ill{} $A^{\circ} \subset B$ \ill{} $A' \subset B'$ \ill{} $A'^{\circ} \subset B'^{\circ}$ \ill{}}
Mais \struck{b') $A'^{\circ} \subset B'$, donc \ill{}} $A'^{\circ} \subset A' \subset A \subset B$ donc $A'^{\circ} \subset B'^{\circ}$ \struck{\ill{} $A \subset B$} \struck{$A^{\circ} \subset B \cap B'$, donc} donc $A'^{\circ} \subset B'^{\circ} \cap B$, donc $B'^{\circ} \cap B \neq \emptyset$, $A'^{\circ} \subset A \subset B$, donc $B' \subset B$, qed.

\textbf{Définition} On dit, sous ces conditions, que $F$ est un \emph{raffinement} de $G$. \uncertain{Toute sous-figure} de $G$ est un raffinement de $G$. \struck{N}otation $F \ll G$. C'est \struck{une} une \emph{relation d'ordre} entre figures\uncertain{,}
\marginal{\textbf{Ex}}

\textbf{Proposition} Soient $F$ et $G$ deux figures, \struck{telles que $F$ est un raffinement de $G$.} \struck{telles} \uncertain{Les} conditions équivalentes :
\begin{itemize}
  \item[a)] \add{($F \ll G$,)} Il n'existe pas de raffinement de $G$ autre que $F$, \uncertain{possédant} \ill{} dont $F$ soit une sous-figure — i.e. $F$ est un élément maximal dans l'ens. des raffinements de $G$ (relation de sous-figure $\leq$)
  \item[b)] \add{($F \ll G$ et)} $|F| = |G|$
\end{itemize}
\struck{c) \ill{}}

\marginal{\textbf{NB} $F$ raffine $G$ \uncertain{ssi} $\forall A \in F$, $F_{A}$ raffine \ill{} \ill{} \ill{} $F$ raffine $G$.}
\note{note oblique dans la marge gauche, à hauteur de la définition ; lecture très incomplète}

\page{23}
b) $\Longrightarrow$ a) Si $F'$ est un raffinement de $G$ tel que $F'$ \uncertain{soit une sur-figure} de $F$, on a $|F| \subset |F'| \subset |G|$, et comme $|F| = |G|$ on a $|F| = |F'|$ \uncertain{d'où} $F = F'$, puisque $F$ est sous-figure de $F'$.

a) $\Longrightarrow$ b) Supposons que $|F| \neq |G|$, \struck{\ill{} $x \in |G| \smallsetminus |F|$, considérons la} je dis qu'il existe une figure $H$ \struck{qui soit} telle que \struck{qui soit un raffinement de $G$,}
\[
  \begin{cases}
    H \neq \emptyset \\
    H \text{ raffine } G \\
    H \text{ disjointe de } F
  \end{cases}
\]
par exemple on prend \struck{un} $x \in |G| \smallsetminus |F|$, $H = \{\{x\}\}$. On a :

\textbf{Corollaire} Soient $F$, $H$ deux figures raffinant $G$, et telles que $F$ et $H$ \uncertain{soient} \uncertain{compatibles}\note{mot écrit au-dessus de « \struck{disjointes} », qui est barré}. Alors $F \cup H$ raffine $G$. \ill{} \uncertain{sup.} (au sens de $\leq$). \struck{Démonstration} \ill{} \ill{}. En effet, on utilise le fait que les multiplicités de $F \cup H$ sont les multiplicités \struck{\ill{}} de $F$ \ill{} $H$. Or la \struck{\ill{}} \ill{} de raffinement se vérifie \ill{} en termes de multiplicités.
\marginal{\ill{} \uncertain{propriété intrinsèque}. \uncertain{Au contraire}, $\pm$ facile. \ill{} \ill{} \ill{} (cas \ill{}) si $F$ et $G$ disjointes…}

\textbf{Définition} On dit dans ce cas que $F$ est une \emph{subdivision} de $G$.
\note{« subdivision » est écrit sous « Définition », souligné ; la ligne du texte laisse la place vide avant « de $G$ »}

\marginal{Donc a) $\Longleftrightarrow$ b) $\Longleftrightarrow$ \ill{} c) \add{$F \ll G$ et} toute strate de $G$ est réunion de strates de $F$ \ill{}. Et \ill{} pour \ill{} $F \ll G$, $G$, \ill{} d) $\begin{cases} 1^{\circ}) \text{ toute strate de } G \text{ est réunion de strates de } F \\ 2^{\circ}) \ \forall A \in F, \ G^{A} = \{B \in G \mid B \supset A\} \text{ a un plus petit élément} \end{cases}$}
\note{longue note oblique dans la marge gauche, en face de la démonstration de a) $\Rightarrow$ b) ; en 2°), un « $\forall A$ » est surchargé}

\page{24}
\textbf{Notation} $F \preccurlyeq G$. C'est encore une relation d'ordre sur $\mathcal{J}$\note{la lettre est lue comme un $J$ cursif, sans doute l'ensemble des figures de $L$ ; lecture incertaine}.

\textbf{NB} Se donner dans $L$ un couple $F \subset G$, avec $F$ raffinant $G$, revient à la même chose que de se donner

\begin{tikzcd}
S\ (= |F|) \arrow[r, hook, "i"] \arrow[d, "p\ \text{surj}"'] & T\ (= |G|) \arrow[r, hook] \arrow[d, "q\ \text{surj}"] & \mathcal{L} \\
I\ (= F) \arrow[r, "\varphi"] & J\ (= G) &
\end{tikzcd}

où $S \subset T$ sont des sous-ens. de $\mathcal{L}$, $\varphi : I \to J$ une application \emph{croissante} d'ensembles ordonnés, et
\[
  p : S \to I, \qquad q : T \to J
\]
des appl. \emph{surjectives} ($\varphi$ pas néc. surjective). Alors l'inclusion $i$ est continue pour les top. \uncertain{canon.} sur $S$, $T$ définies par leurs préordres respectifs (images inv. des top. sur $I$, $J$). L'image inv. d'une partie $G$-\ill{} de $T$ est une partie $F$-\ill{} de $S$, cela \ill{} \ill{} \ill{} : la relation
\[
  J' \mapsto \varphi^{-1}(J')
\]
au niveau des parties de $J$, $I$. L'image inv. d'une partie fermée

\page{25}
(resp. ouverte, resp. loc. fermée) est idem.

Le cas d'une \struck{sub}subdivision est celui où $S = T$, i.e. s'explicite comme une situation

\begin{tikzcd}
S \arrow[r, hook] \arrow[d, "p\ \text{surj}"] & \mathcal{L} \\
I \arrow[d, "\varphi\ \text{surj.}"] & \\
J &
\end{tikzcd}

\uncertain{où} on pose $q = \varphi p$.

\marginal{Attention : même au cas où $\varphi$ est bijective, \uncertain{elle peut ne pas} être un isom. (deux structures d'ordre \ill{} sur un même ens. $I$, avec l'une plus fine que l'autre). C'est le cas d'un $F \preccurlyeq G$, avec $F \neq G$, et revient \ill{} la relation \struck{d'équiv.} $R_{F} = R_{G}$ dans $|F| = |G|$.}

\textbf{Proposition} Soient $F$, $G$ deux figures.
\begin{itemize}
  \item[a)] Pour que $F$ raffine $G$, il faut et il suffit que ce soit une sous-figure d'une subdivision $G'$ de $G$ :
  \[
    F \leq G' \preccurlyeq G
  \]
  \item[b)] \struck{Pour que $F$ soit une subdivision de $G$, il faut et il suffit que ce soit un raffinement, et qu'il soit maximal comme tel.}
\end{itemize}
\note{b) est barré de plusieurs traits obliques ; un trait encadre sa dernière ligne et « Ainsi, dans $F$ », biffé}

\textbf{Dém.} C'est évidemment suffisant. Inversement, supposons $F \ll G$, pour trouver $G'$, on prend $G' = G$ si $F$ est déjà une subd. i.e. $|F| = |G|$. \struck{Sinon}, Dans le cas général, soit
\[
  G' = \bigl\{A \in \mathfrak{P}(\mathcal{L}) \bigm| A \in F \text{ ou } A = \{x\}, \text{ avec } x \in |G| \smallsetminus |F|\bigr\}
  = F \amalg \coprod_{x \in |G| \smallsetminus |F|} \{\{x\}\}
\]

\page{26}
\struck{Alors $F$}

Ainsi, les relations d'ordre $\ll$ et $\preccurlyeq$ dans $\mathcal{J}$ se déterminent mutuellement, en termes de $\leq$. Peut-on de même exprimer $\leq$ en termes de $\ll$ et de $\preccurlyeq$ ?
\marginal{??}

\textbf{NB} On a
\[
  F \leq G \text{ et } F \preccurlyeq G \Longrightarrow F = G .
\]
Il semble que ceci soit vrai dans \struck{\ill{}} contexte — via la condition que si $F \leq G$, $F \neq G$, $\exists H$ \add{non vide}\note{écrit au-dessus d'un mot biffé}, $H$ disjointe de $F$, $H \ll G$ (et nécessairement, si $F \preccurlyeq G$, $H$ serait disjointe de $G$ par C5, d'où contradiction avec ceci : \struck{$H \ll G$,}
\[
  H \ll G, \quad \underset{\text{i.e. } H \text{ raffine } G \text{ et } H \cap G = \emptyset}{H \text{ disjointe de } G} \Longrightarrow H = \emptyset_{F}
\]
\struck{\ill{}} En effet, $H \leq G' \preccurlyeq G$, et $H$ disjointe de $G$ \add{C5} implique $H$ disjointe de $G'$, donc $H = \emptyset$ OK !)

\struck{$F$ \ill{} \ill{}}

Qu'en est-il des autres \ill{} des contextes :
\begin{itemize}
  \item[\textbf{C5}] Si $F \preccurlyeq G$, $H$, alors $H$ disj. de $F$ ssi disj. de $G$ : \uncertain{trivial}, car $|F| = |G|$ et disjonction de $H$ ou $H'$ se voit \uncertain{sur les supports} : $|H| \cap |H'| = \emptyset$.
  \item[\textbf{C6}] \struck{\ill{}}Subdivision induite, on la verra tantôt.
  \item[\textbf{C7}] Ext. d'une subd. OK \ill{} le \ill{} topologique.
  \item[\textbf{C8}] trivial
\end{itemize}

\page{27}
\struck{Proposition} Soient $F$, $G$ deux figures \uncertain{ens.} dans $\mathcal{L}$. \struck{Considérons} Soit \struck{$K$}
\[
  K \struck{I} = \{(A,B) \in F \times G \mid A^{\circ} \cap B^{\circ} \neq \emptyset\}
\]
\struck{Il est clair} que
\[
  |F| \cap |G| = \coprod_{(A,B) \in K} A^{\circ} \cap B^{\circ}
\]
et que les \struck{relations} $A^{\circ} \cap B^{\circ}$ ($(A,B) \in K$) forment les classes de la rel. d'équiv. sur $|F| \cap |G|$ borne sup des relations d'équiv. \struck{induites} \add{induites} par $R_{F}$ et $R_{G}$. Considérons \struck{la famille} \struck{des} les parties
\[
  k = (A,B), \qquad C_{k} = A \cap B
\]
\note{tout ce passage, depuis « Soit $K$ », est barré de trois longs traits obliques}

\begin{tikzcd}
 & S \arrow[rr, hook] \arrow[dd, "\varphi\ \text{surj}"] & & \mathcal{L} \\
 & & S' \arrow[ur, hook] \arrow[d, "\varphi'\ \text{surj}"] & \\
 & I & I' &
\end{tikzcd}

\note{dans le croquis, $S$ et $S'$ sont deux parties de $\mathcal{L}$} $I$, $J$ ens. ord.

Considérons

\begin{tikzcd}
 & S \arrow[rr, hook] \arrow[dd, "\varphi"] & & \mathcal{L} \\
S \cap T \arrow[ur, hook] \arrow[rr, hook] \arrow[dd, "\psi"'] & & T \arrow[ur, hook] \arrow[dd, "\varphi'"] & \\
 & I & & \\
K \arrow[ur] \arrow[rr] & & I' &
\end{tikzcd}

où
\[
  K = \operatorname{Im}\bigl(S \cap T \xrightarrow{(\varphi,\varphi')} I \times J\bigr) = \{(i,i') \in I \times I' \mid A_{i}^{\circ} \cap A'^{\circ}_{i'} \neq \emptyset\}
\]
si $j = (i,i')$, on pose
\[
  \boxed{A_{j} = A_{i} \cap A'_{i'}} \qquad \text{où on pose } A_{i} = \varphi^{-1}(I_{\leq i}), \quad A'_{i'} = \varphi'^{-1}(I'_{\leq i'})
\]
\note{le croquis note $\boxed{K}$ la lettre du coin inférieur gauche ; la formule écrit $I \times J$ où l'on attend $I \times I'$}

\page{28}
\struck{Alors} Ainsi

\textbf{Proposition} On a construit un Inf de $F$ et $G$ pour $\ll$, qu'on appelle la « figure intersection » $F.G$.

Attention, ce n'est pas \add{en général} \struck{la figure} $F \cap G$, \struck{\ill{}} où on rappelle $F, G \subset \mathfrak{P}(\mathcal{L})$, \struck{$G \subset \mathfrak{P}(\mathcal{L})$}. Nous n'allons considérer cette $F \cap G$ que dans le cas où $F$ et $G$ compatibles, mais alors les deux coïncident !

Donc la bonne notion d'intersection est la $F.G$. En fait, se construit pour une famille quelconque de figures, et
\[
  \Bigl|\bigwedge_{i \in I} F_{i}\Bigr| = \bigcap_{i \in I} |F_{i}| .
\]
Cette propriété d'existence d'un Inf pour la relation $\ll$ l'apparente à une propriété de \struck{\ill{}} « \uncertain{demi-treillis} ».

\textbf{Proposition} Si $F' \struck{\ill{}} \preccurlyeq F$ \struck{\ill{}}, alors $F'.G \preccurlyeq F.G$ (\uncertain{Immédiat} sur \ill{}) : Subdivision \emph{induite} par un raffinement.
\note{cette dernière proposition est barrée de plusieurs traits obliques}

\page{29}
\textbf{Proposition} Soient $F'$, $F$, $G$ trois figures.
\begin{itemize}
  \item[a)] $F' \ll F \Longrightarrow F'.G \ll F.G$ \quad (Trivial par déf. de $\ll$ comme Inf)
  \item[b)] $F' \preccurlyeq F \Longrightarrow F'.G \preccurlyeq F.G$ \quad (résulte de a), \ill{} \ill{})
  \item[c)] $F' \leq F \Longrightarrow F'.G \leq F.G$ $^{*}$
\end{itemize}
$^{*}$ \emph{explicitation} $F'.G = \{A \in F.G \mid A \subset |F'|\}$

Prouvons c. Il faut prouver que \struck{$F'.G$} $F'.G \subset F.G$, et est fermé dans $F.G$. \struck{Or} $F'.G$ est formé des $A \cap B$ ($A \in F'$, $B \in G$) tels que $A^{\circ} \cap B^{\circ} \neq \emptyset$, et $F.G$ des $A \cap B$ ($A \in F$, $B \in G$) tels que $A^{\circ} \cap B^{\circ} \neq \emptyset$, et dans les deux cas les \ill{} \ill{} $A^{\circ}$ (\uncertain{pris} \ill{} $F$, \ill{} $F'$) sont les mêmes, donc $F'.G \subset F.G$. Prouvons que c'est fermé i.e. si \struck{$A \cap B$} \struck{$A, A' \in F$,}
\[
  \underset{F'}{A} \cap \underset{G}{B} \in F'.G \quad \text{i.e. } A^{\circ} \cap B^{\circ} \neq \emptyset, \qquad
  \underset{F}{A'} \cap \underset{G}{B'} \in F.G \quad \text{i.e. } A'^{\circ} \cap B'^{\circ} \neq \emptyset,
\]
et $A' \cap B' \subset A \cap B$ \struck{\ill{}}, alors $A' \cap B' \in F'.G$ i.e. $A' \in F'$. Or on sait (d'après la théorie de $F.G$) que $A' \cap B' \subset A \cap B$ $\Longrightarrow A' \subset A$, $B' \subset B$, d'où $A' \in F'$ puisque $F'$ fermé dans $F$.

\textbf{Cor.} L'axiome C6 est valable.

\marginal{\ill{} : $F'.G = G'$ pour un $G'$ tel que $F' \preccurlyeq F$, $G' \preccurlyeq G$, $G' \leq F'$, $G \leq F$ (en carré) mais il faut vérifier un nouvel \ill{} C7${}_{0}$ : cela revient à exiger : a) $P \ll Q \ll R$ et $P \preccurlyeq R \Longrightarrow P \preccurlyeq Q$, $Q \preccurlyeq R$ ; b) $P \ll Q \ll R$ et $P \leq R$, \struck{$\Rightarrow P \leq$} $Q \leq R \Longrightarrow P \leq Q$}
\note{bloc encadré dans la marge droite, relié par un trait à « Vérifions … C7 »}

Vérifions pour de même C7.
\[
  G \leq F, \quad G' \preccurlyeq G \qquad \exists \text{ subdivision } F' \text{ de } F \text{ qui induise } G'.
\]
\note{les trois lettres sont disposées en équerre, $G'$ sous $G$}
On prend $F' = G' \cup \struck{\ill{}} (F \smallsetminus G)$ (c'est une \uncertain{union} disjointe). Il faut vérifier a) c'est une figure b) qu'il \struck{raffine} \add{subdivise} $F$ c) qui induit $G'$ sur $G$.

\marginal{\textbf{Explicitement} : $G \leq F \succcurlyeq F'$, d'où $G'$ avec $G' \preccurlyeq G$, $G' \leq F'$. Alors on a $G' = \{A' \in F' \mid A' \subset |G|\} = \varphi^{-1}(G)$, où $\varphi : F' \to F$ …}
\note{note oblique dans la marge gauche, à hauteur du corollaire}

\page{30}
a) $F'$ est une figure. \struck{i.e.} Notons que \struck{$|F'| =$}
\[
  |F'| = \underbrace{|G'|}_{|G|} \cup \underbrace{|F \smallsetminus G|}_{\supset |F| \smallsetminus |G|} \supset |F|
\]
donc, comme $|F'| \subset |F|$, on a
\[
  |F'| = |F| .
\]
Il faut prouver que $\forall x \in |F'| = |F|$, $\exists$ plus petit élément de $F'$ qui le contienne. Distinguons deux cas
\begin{itemize}
  \item[1°)] $x \in |G| = |G'|$, soit $A \in G'$ tel que $x \in A^{\circ(G')}$. Si $B \in F'$ et \struck{\ill{}} $x \in B$, prouvons $A \subset B$. Si $B \in G'$, c'est évident. Si $B \in F \smallsetminus G$, alors \struck{$B$ contient} $A^{\circ(G')} \cap B \neq \emptyset$ donc $A^{*} \subset B$\note{même signe qu'à la page 21 sur la lettre : astérisque ou rond appuyé} OK
  \item[2°)] $x \in |F| \smallsetminus |G|$, alors les $A \in F'$ qui contiennent $x$ ne sont pas dans $G'$, ils sont donc dans $F \smallsetminus G$. D'autre part, l'unique élément $A$ de $F$ tel que $x \in A^{\circ(F)}$ n'est pas dans $G$ puisque $x \notin |G|$, il est donc dans $F \smallsetminus G$, et \struck{\ill{} $A$} c'est lui (cf. définition de $F \smallsetminus G$) qui \uncertain{est} plus petit él. de $F$ contenant $x$, OK !
\end{itemize}

b) $F'$ \struck{raffine} \add{subdivise} $F$. Comme $|F'| = |F|$, il faut voir \struck{\ill{}} seulement que $F'$ raffine $F$. Donc \struck{que \ill{} que $A \subset B$} \struck{l'ens. $F^{A}$ des $B \in F$ tels que $A \subset B$, $A \in F'$, \ill{} contienne un plus petit} \struck{1) \ill{} un plus petit él. de $F$. C'est évident si $A \in F$ c'est évident. Si \ill{} $F'$ raffine $F$. Si}
\note{trois lignes biffées et encadrées}
\begin{itemize}
  \item[1°)] \struck{$\forall$} $A \in F'$, l'ens. $F^{A} = \{B \in F \mid A \subset B\}$ a un plus petit élément. C'est évident si $A \in F$ \struck{$(B = A)$} (on prend $B = A$) et si $A \in G'$, c'est vrai puisque $G' \ll F$.
  \item[2°)] Si $A \in F'$, $B \in F$, alors $A \cap B$ est \struck{$R_{F'}$-\ill{}} \add{fermé dans $F'$}, i.e. réunion de strates de $F'$. Si $A \in G'$, comme $G' \ll F$, $A \cap B$ est réunion de strates de $G'$, qui sont aussi des strates de $F'$.
\end{itemize}

\page{31}
Si $A \in F \smallsetminus G$, $A \cap B$ est réunion de strates de $F$. \struck{Reste} Reste : voir qu'une strate \add{de $G$} de $F$ est réunion de strates de $F'$. Soit $A$ cette strate de $F$. \struck{On} \struck{$A = \bigcup_{B \in F,\, B \leq F} B^{\circ}$} Si $A \notin G$, donc $A \in F \smallsetminus G \subset F'$, et on gagne. Si $A \in G$, alors (comme $G' \preccurlyeq G$) $A$ est réunion de strates de $G'$, qui sont des strates de $F'$, et on a gagné encore.

\textbf{Corollaire} La \struck{raffinement} \add{subdivision} précédente $F'$ de $F$ est \uncertain{caractérisée} par \add{des} propriétés équivalentes suivantes :
\marginal{\ill{} \ill{}}
\begin{itemize}
  \item[a)] Pour toute subdivision $F''$ de $F$ qui induise $G'$ sur $G$, on a $F'' \preccurlyeq F'$ — i.e. $F'$ est \add{un plus grand élément} \struck{maximal}, pour $\preccurlyeq^{*}$, dans l'ens. des subdivisions de $F$ qui induisent $G'$ sur $G$.
  \note{a) est barré d'une grande croix}
  \item[b)] L'application \uncertain{canonique} associée à $F' \preccurlyeq F$, $\varphi : F' \to F$, induit une bijection
  \[
    F' \smallsetminus G' \xrightarrow{\ \sim\ } F \smallsetminus G .
  \]
  \item[c)] $\forall A \in F \smallsetminus G$, $A^{\circ}$ est une strate \struck{\ill{}} de $F'$ \struck{\ill{}}.
  \item[d)] $F' \supset F \smallsetminus G$, i.e. $\forall A \in F \smallsetminus G$, on a $A \in F'$.
\end{itemize}
\marginal{\textbf{NB.} c'est toujours surj., donc c'est l'injectivité qu'il faut prouver}
\marginal{$^{*}$ ou pour $\ll$, il revient au même pour des subdivisions d'un même $F$}
\marginal{les conditions b), c), d) \uncertain{sont} équivalentes, \ill{} \ill{} \ill{} conditions \ill{} subdivision construite${}^{\#}$. Mais c) très \ill{} (b, c, d) \ill{}}
\note{deux notes obliques dans la marge gauche ; la seconde est coupée d'un trait courbe}

\textbf{Dém.} On va prouver
\[
  \text{d)} \Longrightarrow \text{b)} \Longrightarrow \text{c)} \Longrightarrow \text{a)} \Longrightarrow \text{d)}
\]
\struck{Dém.} d) $\Longrightarrow$ b). Si $(F \smallsetminus G) \subset F'$, alors $\varphi \mid F \smallsetminus G$ est l'identité \add{(car $F \smallsetminus G \subset F$)}, \struck{donc} donc il reste à prouver
\[
  F \smallsetminus G = F' \smallsetminus G' .
\]

\page{32}
\struck{Or} 1°) $F \smallsetminus G \subset F' \smallsetminus G'$ i.e. Tout $A \in F \smallsetminus G$ est $\notin G'$. En effet, si on avait $A \in G'$, on aurait $|A| \subset |G'| = |G|$ donc \struck{$A \in$} (comme $G \leq F$) $A \in G$.

2°) $F' \smallsetminus G' \subset F \smallsetminus G$. Soit $A' \in F' \smallsetminus G'$, prouvons $A' \in F \smallsetminus G$, et d'abord $A' \in F$. \struck{On a} Soit $A$ le plus petit él. de $F^{A'}$ donc $A' \subset A$, $A'^{\circ(F')} \subset A^{\circ(F)}$, prouvons que $A' = A$ (d'où $A' \in F$). \struck{Or}

Je dis que $A \in F \smallsetminus G$ i.e. $A \notin G$, sinon on aurait $A \subset |G| = |G'|$ donc $A' \in G'$ puisque $A' \in F'$ et $G' \leq F'$, absurde. \struck{Mais} Donc $A \in F \smallsetminus G$, donc $A \in F'$, \struck{et} Il en résulte que $A^{\circ(F)} = A^{\circ(F')}$ (car $F' \preccurlyeq F$ \add{i.e. sachant $F' \ll F$} cf lemme plus bas) \uncertain{d'autre part} \struck{$\exists$} $A, A' \in F'$ tels que $A'^{\circ(F')} \subset A^{\circ(F')}$, ce qui prouve $A = A'$.
\marginal{\uncertain{va se simplifier}, cf plus \ill{}}
\note{note verticale dans la marge gauche, séparée du texte par un trait qui court sur toute la hauteur du paragraphe}

Donc on a prouvé que $A' \in F' \smallsetminus G' \Longrightarrow A' \in F$. Prouvons qu'on a $A' \notin G$, d'où $A' \in F \smallsetminus G$ (qed). Si on avait $A' \in G$, on aurait $A' \subset |G| = |G'|$, d'où (comme $A' \in F'$ et $G' \leq F'$) $A' \in G'$, absurde.

\struck{b) $\Longrightarrow$ c) (Soit $A \in F \smallsetminus G$, prouvons que $A^{\circ(F)}$ est une strate \ill{} de $F'$. Soit $A' \in F' \smallsetminus G'$ l'unique élément tel que $\varphi(A') = A$. Prouvons $A'^{\circ(F')} = A^{\circ(F)}$.}
\note{au-dessus, en interligne et biffé avec le reste : « On suppose b) i.e. $\varphi$ induit $F' \smallsetminus G' \xrightarrow{\sim} F \smallsetminus G$ » ; le passage est barré de traits obliques}

b $\Longleftrightarrow$ c \struck{Posons} On a $|F| \smallsetminus |G| = |F'| \smallsetminus |G'|$, soit $U$. Posons $F \smallsetminus G = I^{*}$, $F' \smallsetminus G' = I'^{*}$ (parties ouvertes des ens. d'indices ordonnés $I$ ($= F$), $I'$ ($= F'$)). On a

\page{33}
\begin{tikzcd}
I' \arrow[r, no head, "\supset" description] \arrow[d, "\varphi\ \text{surj}"'] & J' = \varphi^{-1}(J) \arrow[d] \\
I \arrow[r, no head, "\supset" description] & J \text{ partie f.}
\end{tikzcd}

d'où

\begin{tikzcd}
I'^{*} \arrow[d, "\varphi^{*} = \varphi \mid I'^{*}"] \\
I^{*}
\end{tikzcd}

\struck{bijectif.} surjectif
et on a

\begin{tikzcd}
U \arrow[r, "p'"] \arrow[rr, bend right=30, "p"'] & I'^{*} \arrow[r, "\varphi^{*}"] & I^{*}
\end{tikzcd}

\note{les trois flèches sont marquées « surj »}
les \struck{$A$} strates ouvertes de $F'$ qui sont \struck{dans} $\subset U$ (i.e. pour $\in F' \smallsetminus G'$) sont les fibres de $U \xrightarrow{p'} I'^{*}$, les strates ouvertes de $F$ qui sont $\subset U$ (i.e. pour $\notin G$) sont les fibres de $U \xrightarrow{p} I^{*}$.

La condition b) dit que $\varphi^{*}$ est bijective, la condition c) dit que les fibres de $p$ (ce sont les $A^{\circ}$ pour $A \in F \smallsetminus G$) sont des fibres pour $p'$. C'est visiblement équivalent. \struck{D'autre part,}

\marginal{Attention, $\varphi^{*}$ bijectif, $\varphi^{*}$ pas néc. iso !!! Mais on va supposer $\varphi^{*}$ iso, \ill{} \ill{} plus bas.}

D'autre part, que dit d) ? Notons que \add{On a vu \ill{} que pour $F' \ll F$, \struck{\ill{}}} pour $A \in F$ \struck{\ill{}}, on a $A \in F'$ ssi $A^{\circ(F)}$ est une strate ouverte de $F'$ (cf lemme déjà prouvé, qui \uncertain{sera donné} plus bas). Donc d) dit que pour $\forall A_{i} \in F \smallsetminus G = I^{*}$, la $F$-strate $p^{-1}(i) \subset U$ est une $F'$-strate de $U$. Donc cela équivaut visiblement \ill{}.
\note{la fin de la ligne porte un « $p$ \uncertain{bijectif} » biffé}

Donc on a b $\Longleftrightarrow$ c $\Longleftrightarrow$ d.

Montrons que c'est équiv. à a. Soit $F''$ comme dans a), d'où $I''$, et \uncertain{soit} \struck{$S = |F|$}
\[
  S = |F| = |F'| = |F''|
\]

\begin{tikzcd}
S \arrow[r] \arrow[rr, bend left=30, "\varphi"] & I'' \arrow[r, "\varphi''"] & I \\
T \arrow[u, hook] \arrow[r] & J'' \arrow[u, hook] \arrow[r] & J \arrow[u, hook]
\end{tikzcd}

\[
  |G'| = |G| = T
\]
\note{ce dernier passage, depuis « Soit $F''$ », est barré de trois longs traits obliques}

\page{34}
Interprétons la donnée d'un $F' \preccurlyeq F$ induisant $G' \preccurlyeq G$, en termes des données initiales, exprimées par un diagramme d'applications (où $S = |F|$, $T = |G|$, $I$ et $J$ les ens. \add{ord.} d'indices pour $F$, $G$, $J'$ celui pour $G'$)

\begin{tikzcd}
f^{-1}(J) = T \arrow[r, hook] \arrow[d, "g_{J}"] \arrow[dd, bend right=40, "f_{J} = f \mid T"'] & S \arrow[dd, "f\ \text{surj}"] \\
J' \arrow[d, "h_{J}"] & \\
J \arrow[r, hook, "i"] & I
\end{tikzcd}

\note{le diagramme porte en marge gauche la marque « ($*$) » ; sous $J$, de sa main : « partie fermée de $I$ »}

\marginal{\textbf{NB} On part de $f : S \to I$ et \struck{fermé} $J \subset I$, d'où $T = f^{-1}(J)$, $f_{J} : T \to J$, et on factorise $f_{J}$ en $T \to J' \to J$}

La donnée de $G'$ équivaut à celle d'une factorisation de $\varphi' = \varphi \mid T$ en $T \xrightarrow{g_{J}} J' \xrightarrow{h_{J}} J$, où $g_{J}$ est une \struck{\ill{}} appl. croiss. surjective d'ens. ordonnés, et $h_{J}$ d'une application ensembliste surjective. La donnée d'un $F'$ équivaut à une factorisation analogue pour $\varphi$
\[
  S \xrightarrow{\ g\ } I' \xrightarrow{\ h\ } I \qquad g \text{ surj}, \quad h \text{ croiss. surj.}
\]
\struck{Dire} Dire que ça se déduit \struck{\ill{}} \add{induit sur $G$ la donnée $G'$} \struck{\ill{}} signifie qu'il \struck{\ill{}} existe $i' : J' \hookrightarrow I'$ tel que les \struck{\ill{}} diagrammes \struck{\ill{}} suivants commutent

\begin{tikzcd}
T \arrow[r, hook] \arrow[d, "g_{J}"] \arrow[dd, bend right=40, "f_{J}"'] & S \arrow[d, "g"] \\
J' \arrow[r, hook, "i'"] \arrow[d, "h_{J}"] & I' \arrow[d, "h"] \\
J \arrow[r, hook, "i"] & I
\end{tikzcd}

avec $J' = h^{-1}(J)$\note{le $J'$ de cette égalité est écrit en surcharge}.

\page{35}
\struck{Si on a deux \ill{} \ill{} $F'$ et $F''$, correspondant à des ens. d'indices $I'$, $I''$, dire que $F'' \preccurlyeq F'$ signifie que}
\note{ces trois lignes sont barrées de traits obliques et encadrées}

\struck{Donc} Mais dans une telle situation, considérons l'ens. ordonné suivant, défini en termes de la situation ($*$) \struck{\ill{}} « au-dessus de $I$ » : \struck{\ill{}}
\[
  K = J' \amalg I \smallsetminus J \quad \text{comme ens. \ill{}}
\]
la relation d'ordre sur $K$ induisant celles données sur $J'$, $I \smallsetminus J$, et si $\alpha \in J'$, $\beta \in I - J$, on pose
\[
  \begin{cases}
    \alpha \leq \beta \text{ ssi } h_{J}(\alpha) \leq \beta \text{ dans } I \\
    \struck{\ill{}}
  \end{cases}
\]
donc $(\alpha, \beta) \in K \times K$, on a $\alpha \leq \beta$ ssi
\marginal{explicitation pour partir : \ill{}}
\begin{itemize}
  \item[—] ou bien \struck{\ill{}} $(\alpha, \beta) \in J'$ et $\alpha \leq \beta$ dans $J'$,
  \item[—] ou bien $(\alpha, \beta) \in I \smallsetminus J$ et $\alpha \leq \beta$ dans $I \smallsetminus J$
  \item[—] ou bien $\alpha \in J'$, $\beta \in I \smallsetminus J$ et $h_{J}(\alpha) \leq \beta$
\end{itemize}
Ceci posé, on a, en termes de $J' \xrightarrow{h_{J}} J \overset{i}{\hookrightarrow} I$ une factorisation canonique du composé

\begin{tikzcd}
J' \arrow[r, hook, "i"] \arrow[d, "h_{J}"'] & K \arrow[d, "h_{K}"] \\
J \arrow[r, hook, "i"] & I
\end{tikzcd}

$h_{K}$ est $h_{J}$ sur $J'$ et l'ident. can. sur $I \smallsetminus J$.

On a fait ce qu'il fallait \add{sur la relation d'ordre de $K$} pour que sur $J'$, $K$ induise la relation de $J'$, que $h_{K}$ induise un iso d'ens. ordonnés
\[
  K \smallsetminus J' \xrightarrow{\ \sim\ } I \smallsetminus J ,
\]
\struck{et} \add{cela implique que la} \ill{} $J' = h_{K}^{-1}(J)$ (\uncertain{cette condition} \ill{} \struck{\ill{} \ill{}} \ill{}) D'autre part, on voit

\page{36}
immédiat que $K$ est objet final, dans la catégorie des \ill{}, les carrés commutatifs

\begin{tikzcd}
J' \arrow[r, "i'"] \arrow[d, "h_{J}"'] & I' \arrow[d, "h"] \\
J \arrow[r] & I
\end{tikzcd}

\add{d'ens. ordonnés} \struck{\ill{}} (on prend) qui satisfont la condition précédente. Ça devrait impliquer (b, c, d) $\Longleftrightarrow$ a).

\marginal{\textbf{18 juin}}
Dans toutes ces questions, le point de vue des ultra-topologies \add{(i.e. top. \uncertain{préordonnées})} est peut-être le plus pratique, sans prendre la peine des types d'explicitation \uncertain{laborieuse} en termes d'ensembles \add{d'indices} ordonnés. Une « figure » s'explicite comme un couple $(\struck{A}\, S, T)$ d'une partie $S$ de $\mathcal{L}$, et d'une ultra-topologie $T$ sur $S$ (topologie où \emph{toute} réunion de fermés est fermée, i.e. \emph{toute} intersection d'ouverts est ouverte, i.e. l'intersection des voisinages de tout point est un voisinage [ou autrement dit], on présume que tout pt a un plus petit voisinage (nécess. ouvert)). Si $(S,T)$, $(S',T')$ sont comme ci-dessus, dire qu'ils sont compatibles signifie que $T$, $T'$ induisent la même topologie sur $S \cap S'$, qui est fermé dans $S$ (pour $T$) et dans $S'$ (pour $T'$), et alors il y a une unique
\note{« 18 juin » est souligné, en tête du paragraphe, dans la marge gauche}

\page{37}
topologie sur $S \cup S' \simeq S \amalg_{S \cap S'} S'$ qui \struck{induise} induise sur $S$, $S'$ les topologies données, et pour laquelle $S$ et $S'$ soient fermés — c'est la topologie finale pour $S \to S \amalg S'$, $S' \to S \amalg S'$, pour $T$ sur $S$, $T'$ sur $S'$.

[\textbf{NB} Si on a des espaces \struck{ultra}topologiques $(S_{i}, T_{i})_{i \in I}$ et des applications $S_{i} \xrightarrow{f_{i}} X$, la topologie finale sur $S$ (pour $(f_{i})$), admettant comme fermés les $A \in X$ tels que $f_{i}^{-1}(A)$ soit fermé pour tout $i \in I$, est \add{une} ultra-topologie. \struck{La relation de préordre associée est la suivante :} Dans le cas d'espèce, $S \cup S'$ s'identifie à la somme amalgamée de $S$, $S'$ standard en sous-espace fermé commun — opération topologique bien familière !]

Si on a $S' \subset S$, $(S,T)$ et $(S',T')$ sont compatibles i.e. $(S',T')$ est une sous-figure de $(S,T)$, ssi $S'$ est fermé dans $S$ pour $T$, et $T'$ est la top. induite par $T$. Les sous-figures de $(S,T)$ s'identifient donc aux parties fermées de $S$, avec la top. induite.

D'autre part, la relation
\[
  (S',T') \ll (S,T)
\]
signifie que $S' \subset S$, et que $S' \hookrightarrow S$ est continue ($\Longleftrightarrow$ croissante pour les préordres associés : $T'$, $T$). La relation

\page{38}
\[
  (S',T') \preccurlyeq (S,T)
\]
signifie que $S' = S$, et $T'$ plus fine que $T$ i.e. l'application identique $S \xrightarrow{\mathrm{id}} S$ est continue.

\marginal{Description de $\bigwedge_{i}(S_{i},T_{i})$, c'est $S = \bigcap S_{i}$ muni de la \struck{\ill{}} topologie \emph{initiale} sur $S$ pour les $S \hookrightarrow S_{i}$ (munis de $T_{i}$), i.e. \struck{\ill{}} la top. \ill{} \ill{} famille d'applications $f_{i} : S \to S_{i}$ \ill{} des $S_{i}$ munis de $T_{i}$ \ill{} toujours, on a la \emph{topologie engendrée} par les $f_{i}^{-1}(F_{i})$ ($F_{i}$ fermé dans $T_{i}$) : on prend les \struck{réunions} \emph{intersections} \ill{} réunions quelconques \ill{} de la forme $f_{i}^{-1}(F_{i})$ … correspond à la notion de relation de préordre \add{induite} \struck{\ill{}} par une famille de relations de préordres : $x \leq y \Longleftrightarrow f_{i}(x) \leq f_{i}(y)$ $\forall i \in I$}
\note{longue note verticale, écrite de bas en haut le long de la marge gauche sur toute la hauteur de la page ; plusieurs mots en sont biffés ou surchargés}

Revenons à la situation (on reprend les notations $F$ …)
\[
  G \subset F, \qquad G' \preccurlyeq G
\]
donc, posant $T = |G|$, $S = |F|$
\[
  \struck{\ill{}}\ T \hookrightarrow S, \qquad \mathcal{T}_{0} \text{ sur } T, \quad \mathcal{T} \text{ sur } S, \qquad \mathcal{T}'_{0} \leq \mathcal{T}_{0}
\]
\note{sur la page, les deux petits schémas sont disposés en carré, $G'$ sous $G$ et $\mathcal{T}'_{0}$ sous $\mathcal{T}_{0}$ ; ce dernier porte un exposant illisible}
$T$ fermé dans $(S, \mathcal{T})$, $\mathcal{T}_{0}$ top. induite par $\mathcal{T}$, $\mathcal{T}'_{0}$ top. plus fine. On considère les \struck{topologies} $F' \preccurlyeq F$ qui induisent $G'$, i.e. les topologies $\mathcal{T}'$ sur $S$ telles que $\mathcal{T}' \leq \mathcal{T}$ (plus fines que $\mathcal{T}$ — il y a \emph{plus} de fermés) et qui induisent $\mathcal{T}'_{0}$. Donc on veut que 1°) $(S, \mathcal{T}') \to (S, \mathcal{T})$ continue, et 2°) $(T, \mathcal{T}'_{0}) \to (S, \mathcal{T}')$ continue, situation mixte à deux, sur celle des top. initiales et finales …
\[
  \underset{\mathcal{T}'_{0}}{T} \xrightarrow{\ i\ } \underset{\mathcal{T}'\,?}{S} \xrightarrow{\ \mathrm{id}\ } \underset{\mathcal{T}}{S}
\]
\marginal{on sait que le composé est continu.}
Donc on veut (i) que les $A \in \mathfrak{P}(S)$ fermés pour $\mathcal{T}$, \struck{\ill{}} $A \in \mathfrak{P}_{\mathcal{T}}(S)$, soient fermés pour $\mathcal{T}'$, \struck{i.e.} et (ii) que pour \emph{tout} $A'$ fermé pour $\mathcal{T}'$, $i^{-1}(A') = A' \cap T$ soit fermé pour $\mathcal{T}'_{0}$. \add{Si \ill{} \ill{} $\mathcal{T}'_{0} = \mathcal{T}' \mid T$,} \add{$T$ étant fermé pour $\mathcal{T}$} donc pour $\mathcal{T}'$, les parties \struck{fermées} de $T$ \struck{pour $\mathcal{T}'_{0}$} sont fermées
\note{les deux dernières lignes sont surchargées d'additions interlinéaires reliées par des traits ; l'ordre de lecture en est incertain}

\page{39}
pour $\mathcal{T}'$. Dans ces conditions, comme parties fermées \struck{\ill{}} pour $\mathcal{T}'$, on a au moins les
\[
  \begin{cases}
    A \subset S \text{ fermés pour } \mathcal{T} \\
    B \subset T \text{ fermés pour } \mathcal{T}'_{0},
  \end{cases}
\]
donc les réunions $A \cup B$. Je dis que ceux-ci sont \add{déjà} les fermés d'une ultra-topologie \add{$\mathcal{T}'$} qui raffine $\mathcal{T}$, et qui induit la topologie $\mathcal{T}'_{0}$, puisque
\[
  (A \cup B) \cap T = (A \cap T) \cup B
\]
et $A \cap T$ est fermé pour $\mathcal{T}'_{0}$, donc \uncertain{a fortiori} \add{(les fermés dans $(T, \mathcal{T}'_{0})$)}. \struck{\ill{}} C'est une topologie à cause de sa stabilité par réunions finies (et l'ultra-stabilité par intersections \ill{} \ill{} que
\[
  \bigcap_{i \in I} (A_{i} \cup B_{i}) = \bigcup_{J \subset I} \Bigl(\underbrace{\bigcap_{i \in J} A_{i}}_{A}\Bigr) \cap \Bigl(\underbrace{\bigcap_{j \in I \smallsetminus J} B_{j}}_{B}\Bigr)
\]
c'est une réunion de fermés de $(S, \mathcal{T})$, donc un fermé de $(S, \mathcal{T})$, OK.
\note{l'accolade sous le premier facteur est marquée $A$, sous le second $B$}

Quelle est la topologie induite \add{par $\mathcal{T}'$} sur $S \smallsetminus T$ ? $(A \cup B) \cap (S \smallsetminus T) = A \cap (S \smallsetminus T)$, on trouve la top. induite par $\mathcal{T}$.

\struck{Inversement, si $\mathcal{T}' \leq \mathcal{T}$, $\mathcal{T}' \mid T = \mathcal{T}'_{0}$, et $\mathcal{T}' \mid S \smallsetminus T = \mathcal{T} \mid S \smallsetminus T$, je dis que $\mathcal{T}'$ est la top. précédente. Je voudrais voir que ses fermés sont de la forme $A \cup B$ ! Mais soit un tel fermé $C$, alors} \struck{$C = C$}
\note{ce dernier paragraphe est barré de quatre traits obliques}

\marginal{Façon commode d'écrire un fermé $X$ de $(S, \mathcal{T}')$ : \ill{} \ill{} $B = X \cap T$ et $U = X \cap (S \smallsetminus T)$ \ill{} \ill{} \ill{}. Une \uncertain{fermé} des $S \smallsetminus T$, condition $\boxed{B \supset \overline{U} \cap T}$ \add{\uncertain{et}} on a $X = B \cup \overline{U}$. Cette fois correspondance biunivoque entre couples $(B, U)$ satisfaisant les conditions dites, et les $X$. (Dans la \ill{} précédente, on prend $A = \overline{U}$, donc on suppose 1°) $B \supset A \cap T$ 2°) $A \cap (S \smallsetminus T)$ dense dans $A$.}
\note{note oblique qui remplit la marge gauche sur les deux tiers inférieurs de la page}

\page{40}
\struck{car $C$ est $\mathcal{T}'$-fermé,}
\[
  C = \underbrace{(C \cap T)}_{\text{fermé de } (T, \mathcal{T}'_{0}),\ \text{soit } B} \cup \underbrace{\overline{(C \cap (S \smallsetminus T))}^{(\mathcal{T}')}}_{A}
\]
\struck{fermeture dans $S$ \ill{} \ill{} $\mathcal{T}$-fermé, et \ill{} \ill{} $S \smallsetminus T$ \ill{} \ill{} $S \smallsetminus T$. OK. Je voudrais voir ce qui a été vu dans un langage topologique plus général :}
\note{la suite du paragraphe barré de la page 39 ; ce haut de page est barré de deux longs traits obliques}

\textbf{Lemme} Soient $(S, \mathcal{T})$ espace topologique, $T \subset S$ une partie fermée, et $\mathcal{T}'_{0}$ une topologie sur $T$ plus fine que la top. induite (i.e. $(T, \mathcal{T}'_{0}) \to (S, \mathcal{T})$ continue). \struck{Soit $\mathcal{T}'$ une topologie sur} \struck{Il y a alors} Soit $\mathfrak{T}$ l'ens. des topologies sur $S$ qui sont 1°) plus fines que $\mathcal{T}$, et 2°) induisent $\mathcal{T}'_{0}$ sur $T$.
\begin{itemize}
  \item[(a)] Il y a alors dans $\mathfrak{T}$ une topologie la moins fine de toutes, soit $\mathcal{T}'$.
  \item[(b)] Pour qu'une partie $X$ de $S$ soit fermée pour $\mathcal{T}'$, il faut et il suffit que (posant $U = S \smallsetminus T$) :
  \item[1°)] $\struck{X \cap U} \ X \cap U \overset{\text{déf}}{=} X_{U}$ soit fermé dans $U$ (pour $\mathcal{T} \mid U$)
  \item[2°)] $X \cap T \overset{\text{déf}}{=} X_{T}$ soit fermé dans $T$ (pour $\mathcal{T}'_{0}$)
  \item[3°)] $\struck{\ill{}} \ \overline{X_{U}}^{(\mathcal{T})} \cap T \subset X_{T}$.
  \item[(c)] \struck{Les fermés de $\mathcal{T}'$ sont exactement ceux de la forme \ill{} \ill{} \ill{} fermés}
\end{itemize}
\note{la lettre que nous rendons $\mathfrak{T}$ est un $\mathcal{T}$ bouclé et souligné, distinct du $\mathcal{T}$ des topologies}

\end{document}
